\section{A sharper explicit quartic size bound}
\label{sec:kt-quartic-refinement}

The exact stable/endoscopic decomposition improves the coefficient
in KT-QUARTIC-BOUND.  This refinement retains the full stable
rank-three space, including the constructed symmetric-square
subspace, and the full classical rank-two space.  It uses the
already proved stable-only source coverage and the same arithmetic
comparison as the preceding bound.

\begin{theorem}[KT-QUARTIC-QUARTER]
Under the report-wide carried-source convention of
KT-TRACE-MULTIPLICITY and KT-QUARTIC-BOUND, every even
$e\geq2$ satisfies
\[
             \log|K_{2e}(\mathbf Z)|<\frac14e^4.          \tag{KT-Q.1}
\]
More precisely, for every even $e\geq4$ the same argument gives
\[
             \log|K_{2e}(\mathbf Z)|<\frac{239}{960}e^4.
                                                               \tag{KT-Q.2}
\]
This is a new explicit coefficient.  It is not an evaluation
of the existentially selected constant in the archived
VK-EXP7 theorem.
\end{theorem}
\begin{proof}
\paragraph{The comparison integer and its exact degree.}
Set
\[
 N^{\rm st}_{e,7}
  =P^{\rm st}_{e,7}(7^{2e}+7^e+1)\,
                    Q_{e,7}(7^{2e}+1),
\]
where the first polynomial is the full stable compact factor
and the second is the full classical characteristic polynomial
at weight $2e+1$.  The stable-only refinement of
KT-TRACE-MULTIPLICITY and the arithmetic proof of
KT-QUARTIC-BOUND give
\[
             |K_{2e}(\mathbf Z)|\mid(N^{\rm st}_{e,7})^2.
                                                               \tag{KT-Q.3}
\]
The number of complex linear factors in this product is
\[
 s_e=D(e-1)-\left\lfloor\frac{e-1}{3}\right\rfloor
                  +\left\lfloor\frac{2e-2}{3}\right\rfloor.
\]
The dimension formulas give its exact values in the three
even residue classes:
\[
 s_e=\begin{cases}
 e^3/216+e/3,&e\equiv0\pmod6,\\
 e^3/216+4e/9-25/27,&e\equiv2\pmod6,\\
 e^3/216+4e/9-2/27,&e\equiv4\pmod6.
 \end{cases}                                                  \tag{KT-Q.4}
\]
In particular, $s_2=0$ and $s_4=2$.
For $e\geq6$, each line of (KT-Q.4) gives
\[
 \frac{s_e}{e^3}\leq\frac1{216}+\frac4{9e^2}
       \leq\frac1{216}+\frac1{81}
       =\frac{11}{648}<\frac1{32}.
\]
Since $s_4=2=4^3/32$, this proves the uniform inequality
\[
                       s_e\leq\frac{e^3}{32}
                   \qquad(e\geq4\hbox{ even}).                \tag{KT-Q.5}
\]

\paragraph{A common bound for every linear factor.}
The compact operator is a sum of $57$ operators unitary for
an invariant positive-definite Hermitian norm on its
coefficient representation.  Therefore every eigenvalue
$\alpha$ of its scaled operator obeys
$|\alpha|\leq57\,7^{e-1}$.  Restriction to the stable
summand preserves this bound.  The classical bound already
used above is $|a_7|\leq2\,7^e$.  Consequently every linear
factor in the absolute value of $N^{\rm st}_{e,7}$ is at most
\[
 7^{2e}+\frac{64}{7}\,7^e+1
   =7^{2e}\left(1+\frac{64}{7^{e+1}}+\frac1{7^{2e}}\right).
\]
For $e\geq2$ the last parenthesis is bounded by
\[
       1+\frac{64}{343}+\frac1{2401}
       =1+\frac{449}{2401}<\frac43.
\]
The targets also strictly exceed the corresponding eigenvalue
bounds, as in KT-TRACE-MULTIPLICITY, so the comparison integer
is nonzero.  Multiplying the bounds and taking logarithms gives
\[
      \log|N^{\rm st}_{e,7}|
       \leq s_e\left(2e\log7+\log\frac43\right).       \tag{KT-Q.6}
\]
This estimate uses the universal compact operator norm, with
no additional Ramanujan assertion about its stable spectrum.

\paragraph{The numerical coefficient and the two small cases.}
The elementary inequalities
$\log(4/3)<1/3$ and $\log7<39/20$ suffice.  The second can
be checked without decimal approximations: the positive
exponential series gives
\[
 \exp(39/20)>\sum_{j=0}^{8}\frac{(39/20)^j}{j!}
    =7+\frac{3130848254169}{114688000000000}>7.
\]
For every even $e\geq4$, combine (KT-Q.3), (KT-Q.5), and
(KT-Q.6), and use $1/(3e)\leq1/12$:
\[
 \begin{aligned}
 \log|K_{2e}(\mathbf Z)|
 &\leq2s_e\left(2e\log7+\log\frac43\right)\\
 &<\frac{e^4}{16}\left(\frac{39}{10}+\frac1{3e}\right)\\
 &\leq\frac{e^4}{16}\left(\frac{39}{10}+\frac1{12}\right)
   =\frac{239}{960}e^4<\frac14e^4.
 \end{aligned}
\]
This includes $e=4$, whose exact degree was treated above.
For $e=2$, both polynomial spaces are empty, so
$N^{\rm st}_{2,7}=1$.  Equation (KT-Q.3) gives
$|K_4(\mathbf Z)|=1$, and its logarithm is zero, proving
(KT-Q.1) in this remaining case.
\end{proof}

\paragraph{Dependencies retained.}
The stable-only source injection is the refinement following
KT-U3-SPECTRAL-SPLIT.  The arithmetic divisibility in
(KT-Q.3) retains the closed OneChar extraction, the all-even-degree
GD statement, the inherited exceptional-prime argument, the
carried small-prime Vandiver input, and the separate
Rognes--Weibel two-primary input already stated in
KT-QUARTIC-BOUND.  This section changes the dimension and
archimedean estimates; it adds no new arithmetic vanishing
claim and does not omit the stable symmetric-square branch.
